\documentclass{article}

\usepackage{amsmath,amssymb}
\usepackage{xurl}
\usepackage[hidelinks]{hyperref}

% Shared commands for notes in docs/paper-gaps/.

\DeclareUnicodeCharacter{2080}{\ifmmode{}_{0}\else\textsubscript{0}\fi}
\DeclareUnicodeCharacter{2081}{\ifmmode{}_{1}\else\textsubscript{1}\fi}
\DeclareUnicodeCharacter{2082}{\ifmmode{}_{2}\else\textsubscript{2}\fi}
\DeclareUnicodeCharacter{2099}{\ifmmode{}_{n}\else\textsubscript{n}\fi}

\newcommand{\Lean}{\textsc{Lean}}
\ExplSyntaxOn
\NewDocumentCommand{\leanid}{m}{
  \ifmmode
    \textsf{\detokenize{#1}}
  \else
    \group_begin:
    \urlstyle{sf}
    \tl_set:Nn \l_tmpa_tl {#1}
    \tl_replace_all:Nnn \l_tmpa_tl {\_} {_}
    \exp_args:NV \path \l_tmpa_tl
    \group_end:
  \fi
}
\ExplSyntaxOff
\providecommand{\tr}{\operatorname{tr}}
\newcommand{\tnleanIssueBase}{https://github.com/LionSR/TNLean/issues/}
\newcommand{\tnleanPrBase}{https://github.com/LionSR/TNLean/pull/}
\newcommand{\tnleanDocBase}{https://sirui-lu.com/TNLean/docs/}
\newcommand{\ghissue}[1]{\href{\tnleanIssueBase#1}{GitHub issue \##1}}
\newcommand{\ghpr}[1]{\href{\tnleanPrBase#1}{PR \##1}}
\newcommand{\doclink}[1]{\href{\tnleanDocBase#1.html}{\path{#1}}}

% Dirac notation and the identity operator, as in blueprint/src/macros/common.tex.
% \providecommand keeps note-local definitions authoritative.
\usepackage{dsfont}
\providecommand{\ket}[1]{|#1\rangle}
\providecommand{\bra}[1]{\langle #1|}
\providecommand{\braket}[2]{\langle #1 | #2 \rangle}
\providecommand{\ketbra}[2]{|#1\rangle\!\langle #2|}
\providecommand{\Id}{\mathds{1}}
\providecommand{\one}{\mathds{1}}
\providecommand{\C}{\mathbb{C}}
\providecommand{\MN}[1]{M_{#1}(\C)}
\providecommand{\MD}{M_D(\C)}

% External referencing for paper sources.  The build helper writes sanitized
% auxiliary files under build/paper-aux/ and excludes duplicated source labels.
\usepackage{xr}

% Import a generated paper auxiliary file with a caller-supplied label prefix.
% The candidate paths support builds from docs/paper-gaps/, the repository root,
% and build/paper-aux/ itself.
\newcommand{\importpaperaux}[2]{%
  \IfFileExists{../../build/paper-aux/#2.aux}{%
    \externaldocument[#1]{../../build/paper-aux/#2}%
  }{%
    \IfFileExists{build/paper-aux/#2.aux}{%
      \externaldocument[#1]{build/paper-aux/#2}%
    }{%
      \IfFileExists{#2.aux}{%
        \externaldocument[#1]{#2}%
      }{}%
    }%
  }%
}

% General external reference macros
\newcommand{\paperref}[2]{\ref{#1-#2}}
\newcommand{\papereqref}[2]{(\ref{#1-#2})}

% CPSV16 source (1606.00608 / MPDO-22-12-17-2.tex) external document and shortcuts
\importpaperaux{cpsv16-}{1606.00608/MPDO-22-12-17-2-tnlean}
\newcommand{\cpsvref}[1]{\paperref{cpsv16}{#1}}
\newcommand{\cpsveqref}[1]{\papereqref{cpsv16}{#1}}

% Verdict marker: \gapnote{<kind>}{<status>}. Kind is the policy
% classification (clarification, local-correction, scope-restriction,
% unfaithful, false-source, open-gap); status is open, wip (elimination
% underway), resolved, or historical. Machine-read by the site index;
% prints a verdict line.
\newcommand{\gapnote}[2]{\par\noindent\textbf{Verdict:} #1 (#2).\par}


\title{The Pure-State Blocking-up-to-Gauge Equivalence Is False}
\author{}
\date{2026-08-10}

\begin{document}
\maketitle
\gapnote{false-source}{resolved}

\section{Source claim}

Appendix~D of Ref.~\cite{Cirac2016MPDO_arXiv}, in the source equation
\path{RFP-gauge} and local lines~2100--2107, allows a renormalization step to
introduce a gauge transformation on the virtual indices. The paragraph after
the corresponding diagram claims that, for pure tensors in canonical form,
this blocking-up-to-gauge condition is equivalent to the ordinary
renormalization fixed-point condition. Its first step asserts that similarity
of the transfer maps $\mathcal E^2$ and $\mathcal E$ forces their eigenvalues
to lie in $\{0,1\}$.

That inference is false. A finite spectrum can be invariant under squaring
because its nonzero eigenvalues form a nontrivial root-of-unity cycle.
Similarity of $\mathcal E^2$ and $\mathcal E$ therefore does not imply
\begin{align}
    \mathcal E^2 &= \mathcal E.
    \label{eq:cpsv16_rfp_gauge_pure_equivalence_false_unjustified_idempotence}
\end{align}

Definition \path{RFPMixedTS} at source lines~657--660 of
Ref.~\cite{Cirac2016MPDO_arXiv} imposes canonical form on the MPDO tensor $M$.
The formal refutation uses this literal reading of ``in CF''. It proves the
canonical-form and weight-normalization hypotheses for the doubled MPO tensor
$M=A\otimes\overline A$, not merely for the underlying pure MPS tensor $A$.

\section{Source-shaped counterexample}

Let
\begin{align}
    \omega
        &:= e^{2\pi i/3},
    \label{eq:cpsv16_rfp_gauge_pure_equivalence_false_cube_root}\\
    A
        &:= \begin{pmatrix}
            1 & 0\\
            0 & \omega
        \end{pmatrix},
    \label{eq:cpsv16_rfp_gauge_pure_equivalence_false_phase_tensor}\\
    S
        &:= \begin{pmatrix}
            0 & 1\\
            1 & 0
        \end{pmatrix}.
    \label{eq:cpsv16_rfp_gauge_pure_equivalence_false_swap_matrix}
\end{align}
Since $\omega^3=1$, direct multiplication gives
\begin{align}
    A
        &= \omega S A^2S^{-1}.
    \label{eq:cpsv16_rfp_gauge_pure_equivalence_false_block_gauge_identity}
\end{align}
Equation~\ref{eq:cpsv16_rfp_gauge_pure_equivalence_false_block_gauge_identity}
has the orientation of the source diagram: the original tensor is the
phase-scaled virtual conjugate of its two-site blocking.

Define the corresponding pure MPDO tensor by
\begin{align}
    M
        &:= A\otimes\overline A.
    \label{eq:cpsv16_rfp_gauge_pure_equivalence_false_doubled_tensor}
\end{align}
For $N\geq 1$, let $V^{(N)}(A)$ denote the MPS wavefunction generated by $A$.
Then $M$ generates
\begin{align}
    \ket{V^{(N)}(A)}\bra{V^{(N)}(A)},
    \label{eq:cpsv16_rfp_gauge_pure_equivalence_false_rank_one_mpdo}
\end{align}
which is positive semidefinite. Thus $M$ satisfies the MPDO standing
hypothesis in Definition \path{RFPMixedTS} of
Ref.~\cite{Cirac2016MPDO_arXiv}.

To verify canonical form, set
\begin{align}
    \mu_0 &:= 1,
    \label{eq:cpsv16_rfp_gauge_pure_equivalence_false_first_weight}\\
    \mu_1 &:= \omega,
    \label{eq:cpsv16_rfp_gauge_pure_equivalence_false_second_weight}
\end{align}
and let $I_1$ denote the bond-one scalar unit tensor. After identifying each
doubled bond coordinate with a pair $(p,q)\in\{0,1\}^2$, one obtains
\begin{align}
    M
        &=
        \bigoplus_{p,q\in\{0,1\}}
        \mu_p\overline{\mu_q}\,I_1.
    \label{eq:cpsv16_rfp_gauge_pure_equivalence_false_canonical_decomposition}
\end{align}
Each copy of $I_1$ is normal because its transfer map is the identity, and
\begin{align}
    \left|\mu_p\overline{\mu_q}\right|
        &= 1.
    \label{eq:cpsv16_rfp_gauge_pure_equivalence_false_weight_normalization}
\end{align}
Equation~\ref{eq:cpsv16_rfp_gauge_pure_equivalence_false_canonical_decomposition}
therefore supplies canonical-form data whose four weights satisfy the
normalization convention in Ref.~\cite{Cirac2016MPDO_arXiv}.

For one physical letter, the physical matrix algebras before and after
blocking are one-dimensional. Every unitary-conjugation channel on either
algebra is the identity, so the source's physical channels cannot contribute a
phase. The MPS amplitude phase instead cancels after doubling:
\begin{align}
    &\left(\omega S A^{[2]}S^{-1}\right)
    \otimes
    \left(\overline\omega
        \overline S\,\overline{A^{[2]}} \overline S^{-1}\right)
    \notag\\
    &\qquad=
    \left(\omega\overline\omega\right)
    (S\otimes\overline S)
    M^{[2]}
    (S\otimes\overline S)^{-1}
    \label{eq:cpsv16_rfp_gauge_pure_equivalence_false_phase_cancellation}\\
    &\qquad=
    (S\otimes\overline S)
    M^{[2]}
    (S\otimes\overline S)^{-1}.
    \label{eq:cpsv16_rfp_gauge_pure_equivalence_false_doubled_gauge_identity}
\end{align}
Thus $M^{[2]}$ and $M$ satisfy the literal virtual-gauge condition in both
directions of the Appendix~D diagram in
Ref.~\cite{Cirac2016MPDO_arXiv}. The MPS identity
Eq.~\ref{eq:cpsv16_rfp_gauge_pure_equivalence_false_block_gauge_identity} is
used only to construct this exact doubled-MPDO witness.

For a bond matrix $\rho$, define the one-letter transfer map by
\begin{align}
    \mathcal E_A(\rho)
        &:= A\rho A^\dagger.
    \label{eq:cpsv16_rfp_gauge_pure_equivalence_false_transfer_definition}
\end{align}
Let $e_{21}$ be the off-diagonal matrix unit. Then
\begin{align}
    \mathcal E_A(e_{21})
        &= \omega e_{21},
    \label{eq:cpsv16_rfp_gauge_pure_equivalence_false_first_transfer_action}\\
    \mathcal E_A^2(e_{21})
        &= \omega^2 e_{21}.
    \label{eq:cpsv16_rfp_gauge_pure_equivalence_false_second_transfer_action}
\end{align}
Because $\omega\ne 1$, the two values differ, so $\mathcal E_A$ is not
idempotent. Equivalently,
\begin{align}
    \mathcal E_A^{2^n}(e_{21})
        &=
        \omega^{2^n}e_{21}
    \label{eq:cpsv16_rfp_gauge_pure_equivalence_false_dyadic_oscillation}
\end{align}
oscillates between the two nontrivial cube-root phases and does not converge.
An idempotent transfer map, by contrast, has every positive power equal to
itself.

The formal declarations are listed in the accompanying formal
development.\footnote{Formal declarations:
\leanid{MPOTensor.IsOneLetterRFPViaTSUpToVirtualGauge},
\leanid{MPSTensor.IsPureOneLetterRFPViaTSUpToVirtualGauge},
\leanid{MPSTensor.IsBlockedGaugePhaseFixedPoint},
\leanid{MPSTensor.doubledVirtualGauge},
\leanid{MPSTensor.isPureOneLetterRFPViaTSUpToVirtualGauge_of_isBlockedGaugePhaseFixedPoint},
\leanid{MPSTensor.cube_phase_blocked_gauge_identity},
\leanid{MPSTensor.cubePhaseDoubledTensor_isMPDO},
\leanid{MPSTensor.cubePhaseDoubledTensor_isCPSVCanonicalForm},
\leanid{MPSTensor.cubePhaseDoubledTensor_exists_weightNormalized},
\leanid{MPSTensor.cubePhaseTensor_isPureOneLetterRFPViaTSUpToVirtualGauge},
\leanid{MPSTensor.cubePhaseTensor_not_isTransferIdempotent}, and
\leanid{MPSTensor.cpsv16_pure_rfp_gauge_equivalence_false} in
\doclink{TNLean/MPS/RFP/GaugeBlockingCounterexample}.}

\section{Verdict}

The pure-state equivalence following the source equation \path{RFP-gauge} in
Ref.~\cite{Cirac2016MPDO_arXiv} is false as printed. It fails even when the
one-letter doubled MPO tensor generates an MPDO family, is in canonical form,
and satisfies the normalization stated at source line~246.

The valid gauge-covariance statement remains unaffected: a virtual gauge
preserves an already established renormalization fixed point. Gauge covariance
does not turn a blocking-up-to-gauge hypothesis into transfer-map
idempotence.

The printed pure-state spectral argument therefore cannot support the
mixed-state continuation in the next paragraph of
Ref.~\cite{Cirac2016MPDO_arXiv}. This counterexample does not decide the
mixed-state question.

\bibliographystyle{alpha}
\bibliography{references}

\end{document}
